\documentclass[10pt,a4paper]{amsart}
\usepackage[margin=19mm]{geometry}
\usepackage{amsmath,amssymb,mathtools}
\usepackage[scr=rsfs,cal=euler]{mathalfa}
\usepackage{xfrac}
\usepackage[unicode,hidelinks,bookmarks=false]{hyperref}
\newcommand{\ie}{i.e.\ }
\newcommand{\cf}{cf.\ }
\newcommand{\resp}{respectively\ }
\newcommand{\Subsection}{\subsection}
\providecommand{\nobreakdash}{\nobreak\hbox{-}}
\setlength{\emergencystretch}{2em}
\multlinegap0pt
\usepackage{mathtools}
\usepackage{amssymb}
\usepackage{xcolor}
\usepackage{algorithm}\usepackage{algpseudocode}
\usepackage{bm}
\usepackage{float}
\usepackage{cancel}
\usepackage{csquotes}
\newtheorem{lemma}{Lemma}
\usepackage{cleveref}
\crefname{lemma}{Lemma}{Lemmas}
\newcommand{\R}{\ensuremath{\mathbb{R}}}% real numbers
\title{Sparse Randomised Approximation of Normal Cycles}
\author{Allen Paul, Neill Campbell, Tony Shardlow}
\date{Source: arXiv:2503.01057v3 (CC BY 4.0)}
\begin{document}
\maketitle

\noindent\emph{Source and licence.} Sparse Randomised Approximation of Normal Cycles. Version arXiv:2503.01057v3 (CC BY 4.0), distributed by arXiv under the Creative Commons Attribution 4.0 International licence, http://creativecommons.org/licenses/by/4.0/, as recorded in the arXiv licence metadata for this exact version and shown on the abstract page https://arxiv.org/abs/2503.01057v3. Full licence notice: \texttt{licenses/arxiv-2503.01057-CC-BY-4.0.txt}.\par\smallskip

\noindent\textit{Editorial scope.} Counted unit: the article's lemma lemma (source0.tex lines 1032--1041). The article's complete original proof of it is delivered with it (source0.tex lines 1043--1076, one \texttt{proof} environment of 34 lines). No mathematics is written, repaired or re-derived: the statement and the proof are the article's own lines, in the article's own notation, and no retained line is substituted or re-flowed.  Retained uncounted as the source-local context the proof needs: lines 977--979.  Nothing was omitted from any retained span, so the package carries no omission record.  No retained line contains a citation, so the wrapper carries no bibliography and no bibliography entry.  No retained line contains a figure environment, an image-inclusion command or drawing code, so no image is delivered and the package declares no asset dependency.  Editorial formatting only, in addition: an AMS \texttt{amsart} preamble that restates the article's own declarations and macros used by the retained lines, namely the environments \texttt{lemma} on separate counters, together with the standard packages those lines need.  The retained environments print in this document as Lemma 1 in order of appearance, whereas the article prints the same items with the numbering of its own full document.  The complete native source of the article as distributed in the arXiv e-print for this exact version is bundled unchanged as \texttt{sources/174118/source0.tex}.
\begin{align}\label{exteriorinnerproddef}
    \langle v_{1}\wedge \dots \wedge v_{k},w_{1}\wedge \dots \wedge w_{k}\rangle_{(\Lambda^{k}\R^{d}) } \coloneq \det((\langle v_{i} ,w_{j}\rangle)_{1\leq i,j\leq k})
\end{align}

\begin{lemma}\label[lemma]{simplelemmma}
Let $S \subset \mathbb{R}^d$ be an oriented $m$-rectifiable set with associated current $[S] \in \Omega_0^m(\mathbb{R}^d)^*$ defined by
\[
[S](\omega) \coloneqq \int_S \bigl(\omega(x) \mid \tau_S(x)\bigr)\, d\mathcal{H}^m(x), \qquad \omega \in \Omega_0^m(\mathbb{R}^d),
\]
where $\tau_S(x) = e_1(x) \wedge \cdots \wedge e_m(x)$ is the wedge product of a positively oriented orthonormal basis of the tangent space at $x$. Then,
\[
\bigl|[S](\omega)\bigr| \leq \|\omega\|_\infty\, \mathcal{H}^m(S), \qquad \omega \in \Omega_0^m(\mathbb{R}^d).
\]
\end{lemma}

\begin{proof}
By the standard estimate $\bigl|\int f\bigr| \leq \int |f|$ (triangle inequality for integrals), we have
\[
\bigl|[S](\omega)\bigr| 
= \left|\int_S \bigl(\omega(x) \mid \tau_S(x)\bigr)\, d\mathcal{H}^m(x)\right|
\leq \int_S \bigl|\bigl(\omega(x) \mid \tau_S(x)\bigr)\bigr|\, d\mathcal{H}^m(x).
\]

The pairing $\bigl(\omega(x) \mid \tau_S(x)\bigr)$ is the duality pairing between $\omega(x) \in (\Lambda^m \mathbb{R}^d)^*$ and $\tau_S(x) \in \Lambda^m \mathbb{R}^d$. By definition of the dual norm on $(\Lambda^m \mathbb{R}^d)^*$,
\[
\bigl|\bigl(\omega(x) \mid \tau_S(x)\bigr)\bigr| 
\leq |\omega(x)|_{(\Lambda^m \mathbb{R}^d)^*}\, |\tau_S(x)|_{\Lambda^m \mathbb{R}^d}.
\]
Since $\{e_1(x), \ldots, e_m(x)\}$ is an orthonormal basis of the tangent space, the inner product on $\Lambda^m \mathbb{R}^d$ (cf. \cref{exteriorinnerproddef}) gives
\begin{align*}
|\tau_S(x)|^2 
&= \langle e_1(x) \wedge \cdots \wedge e_m(x),\; e_1(x) \wedge \cdots \wedge e_m(x) \rangle_{\Lambda^m \mathbb{R}^d}\\
&= \det\!\bigl(\langle e_i(x), e_j(x)\rangle\bigr)
= \det(I_m) = 1.
\end{align*}
Therefore, $|\tau_S(x)|_{\Lambda^m \mathbb{R}^d} = 1$, and the pointwise bound simplifies to
\[
\bigl|\bigl(\omega(x) \mid \tau_S(x)\bigr)\bigr| \leq |\omega(x)|_{(\Lambda^m \mathbb{R}^d)^*}.
\]

Substituting the pointwise estimate into the integral,
\[
\bigl|[S](\omega)\bigr| 
\leq \int_S |\omega(x)|_{(\Lambda^m \mathbb{R}^d)^*}\, d\mathcal{H}^m(x)
\leq \sup_{x \in \mathbb{R}^d} |\omega(x)|_{(\Lambda^m \mathbb{R}^d)^*} \cdot \mathcal{H}^m(S)
= \|\omega\|_\infty\, \mathcal{H}^m(S).\qedhere
\]

\end{proof}

\end{document}
